\section{Parity witnesses, weighted repair budgets, and their exact limit}
\label{sec:parity-repair}

The odd-block argument extends to arbitrary integer witnesses, including
overlapping or unequal odd blocks. The extension supplies explicit
constraints on \emph{which} rows a repair must remove. Saturated-lattice
duality also identifies exactly what these certificates detect: absence
of an all-odd integer kernel vector. A certified order-twenty-four
example shows why that arithmetic test still falls short of sign
extension. Arithmetic completion is a classical subject; a relevant
primary source is Verheiden's thesis,
\emph{Arithmetical Properties of Combinatorial Matrices} (1978),
\url{https://thesis.caltech.edu/2899/}, and his article
\emph{Integral and rational completions of Hadamard matrices},
\url{https://doi.org/10.1016/0097-3165(80)90028-X}.
The results below have self-contained proofs; no novelty claim is made.

\begin{theorem}[Weighted odd-witness budget]
Let $P$ have $r$ orthogonal sign rows $p_1,\ldots,p_r$ of length $N$,
and let $f\in\mathbb Z^N$ have odd coordinate sum. Set
\[
 R=NI_N-P^{\mathsf T}P,\qquad
 C_f=f^{\mathsf T}Rf
     =N\|f\|^2-\sum_i(p_i\cdot f)^2\ge0.
\]
If rows indexed by $D$ are deleted and $L$ mutually orthogonal sign
rows are inserted, orthogonal also to all retained rows, then
\[
 L\le C_f+\sum_{i\in D}(p_i\cdot f)^2.
\]
Consequently full completion requires
\[
 \sum_{i\in D}\bigl((p_i\cdot f)^2-1\bigr)
       \ge N-r-C_f.
\]
Every odd-sum integer $f$ gives such an inequality, so a collection
of witnesses gives simultaneous weighted constraints on row selection.
\end{theorem}
\begin{proof}
Bessel's inequality for orthogonal rows of squared norm $N$ bounds
the sum of squared inner products with $f$ by $N\|f\|^2$.
For every sign row $v$, $v\cdot f\equiv\sum_j f_j\equiv1\pmod2$,
so $(v\cdot f)^2\ge1$. Subtracting the contributions of the retained
rows proves the first inequality. A full completion inserts
$L=N-r+|D|$, giving the second.
\end{proof}

For a signed indicator $f$ of an odd block of size $s$, the constant
is $C_f=Ns-\sum_i(p_i\cdot f)^2$. The theorem applies independently
to every block, with no equality-of-sizes or disjointness assumption.
If $f$ lies in the rational row span of $P$, then $C_f=0$ and no
additional sign row exists. The budget then depends only on the
deleted rows' squared overlaps with the witness.

\begin{proposition}[The first possible parity deficit is eight]
Assume $4\mid N$ and write $d=N-r$. For every odd-sum integer $f$,
\[
 C_f\equiv d\pmod8.
\]
Thus for $0\le d\le7$ these witness inequalities automatically give
$C_f\ge d$; they cannot certify failure of full completion without
row deletion. At $d=8$, the value $C_f=0$ is possible.
\end{proposition}
\begin{proof}
Every $p_i\cdot f$ is odd, hence its square is $1$ modulo $8$.
Also $\|f\|^2$ is odd, so $N\|f\|^2\equiv N\pmod8$ because
$4\mid N$. Therefore $C_f\equiv N-r\pmod8$.
Nonnegativity gives the remaining assertion. The repeated $4\times12$
seed supplies the value zero at defect eight. This congruence alone
does not prove the classical seven-row completion theorem.
\end{proof}

\subsection{What the certificate detects exactly}
\begin{theorem}[Saturated-lattice alternative]
For any rational row space $V\subseteq\mathbb Q^N$, exactly one of
the following holds:
\begin{enumerate}
\item $V\cap\mathbb Z^N$ contains a vector with odd coordinate sum;
\item $V^\perp\cap\mathbb Z^N$ contains a vector all of whose
coordinates are odd.
\end{enumerate}
For a partial Hadamard row space, the first case supplies a zero-budget
witness. The second case supplies arithmetic parity feasibility,
but does not assert a sign vector exists.
\end{theorem}
\begin{proof}
Put $\Lambda=V\cap\mathbb Z^N$ and
$K=V^\perp\cap\mathbb Z^N$. Both are primitive: an integer multiple
lying in the relevant rational subspace implies the original vector
lies there. A primitive subgroup of $\mathbb Z^N$ has a basis
extendible to an integer unimodular basis, so reduction modulo $2$
preserves its rank. If $\dim V=r$, the images $\bar\Lambda,\bar K$
in $\mathbb F_2^N$ have dimensions $r,N-r$ and are orthogonal.
Thus $\bar K=\bar\Lambda^\perp$. The all-ones vector lies in
$\bar K$ exactly when every vector of $\bar\Lambda$ has even
coordinate sum. This is precisely the stated alternative.
\end{proof}

There is a simple exact algorithm for this alternative. Starting with
an independent integer row basis, find a dependence modulo $2$.
The corresponding sum of basis rows has even entries; divide it by
$2$ and replace any participating row. Rational span is preserved,
and the squared covolume decreases by a factor $4$. As it stays a
positive integer, the procedure terminates with independent reductions
modulo $2$. Its row lattice has odd index in the saturated lattice,
so it has the same reduction modulo $2$. An odd-sum basis row is
therefore a complete certificate for the first case. Otherwise select
a full-rank odd minor, set free kernel coordinates to $1$, and solve
for the pivot coordinates. Their denominators are odd, and clearing
those denominators produces the all-odd integer kernel vector.
The accompanying checker implements this procedure from scratch.

\subsection{Sparse certificates beyond uniform repeated seeds}
For the earlier terminal $16\times24$ frame, encoded by row masks
\[
\begin{gathered}
0,16773120,16519104,10359454,7684730,5820908,9389414,1547504,\\
6736278,11838780,3701846,13790378,10646512,2963882,12892774,5974836,
\end{gathered}
\]
with bit $j$ indicating entry $-1$ at zero-based coordinate $j$,
there are two particularly short witnesses:
\[
 f=e_2+e_8+e_{19},\qquad g=e_9-e_{10}+e_{22}.
\]
Direct integer calculation gives $C_f=C_g=0$ and
$\|f\|^2=\|g\|^2=3$. This certifies unextendibility without
enumerating candidate sign rows. The row indices having absolute
overlap $3$ are respectively
\[
 A=\{0,4,5,6,7,11,15\},\qquad
 B=\{4,6,7,9,10,14,15\};
\]
all other overlaps have magnitude $1$. Any augmenting exchange must
delete a row from each of $A$ and $B$. These necessary conditions
do not by themselves establish the previously certified three-row
exchange barrier, but they give concrete filters for a repair search.

There is also a positive case where the weighted budget selects an
optimal constructive route. In the repeated order-twelve tensor
family after $k$ fibers have been repaired, a triple-indicator
witness has zero budget. Original repeated rows have overlap
magnitude $3$, while installed affine-plane rows have overlap
magnitude $1$. Deleting $a$ original rows and $b$ installed rows
can therefore insert at most $9a+b$ rows. Net growth is at most
$8a$, so deleting only installed rows cannot augment the frame.
Full completion requires at least $t-k$ more original-row deletions,
regardless of how many installed rows are removed. The explicit fiber
repair attains this bound: sacrifice one original row and insert the
nine words $(x,y,x+y,x-y)$ over $\mathbb F_3$, interpreted as one
negative position per triple. This supplies a certified choice of
which class of rows to preserve during constructive repair.

\subsection{Odd integer feasibility still misses a Euclidean gap}
The lattice alternative has an explicit limitation. Consider the
$12\times24$ partial frame with row masks
\[
\begin{gathered}
0,16773120,16519104,9656012,13838682,14783914,\\
1944718,11409518,15230494,5921970,3487346,13018356.
\end{gathered}
\]
Its rows are orthogonal. Its integer kernel contains
\[
\begin{split}
v={}&(-1,-1,3,-1,1,1,-3,-1,1,1,-1,1,\\
&1,-1,-1,-1,1,-1,1,-1,1,-1,1,1),
\end{split}
\]
with $\|v\|^2=40$. Nevertheless its odd integer kernel has no vector
of squared norm $24$ or $32$, so its shortest odd vector has squared
norm exactly $40$.

Here this last assertion is a finite exact certificate. Rational
elimination selects twelve free coordinates. Enumerating their
$2^{11}=2048$ sign assignments up to global negation rules out norm
$24$. A norm-$32$ all-odd vector would have exactly one coordinate
of magnitude $3$ and every other coordinate of magnitude $1$.
If that exceptional coordinate is a pivot, it is covered by the same
2048 assignments. If it is free, there are another
$12\cdot2^{11}=24576$ assignments to check. None gives norm $32$.
Every all-odd vector has squared norm congruent to $24$ modulo $8$,
and the displayed vector attains the next possibility, $40$.
The verifier performs all these checks with exact integer and rational
arithmetic. Thus the frame has no sign extension despite passing the
all-odd kernel feasibility test. Real progress beyond parity requires
controlling the short-vector problem, not merely clearing denominator
or congruence obstructions.
